\section{ADR-008: DAW integration and engine discovery}
\textbf{Status:} accepted (M11).

\textbf{Context:} producers want stems inside their \gls{daw}. A REAPER script
and a VST3 plugin run in the DAW's process, not in the desktop shell, so they do
not receive the port and token that the shell hands to its WebView (ADR-004).
They need a way to find the running engine that keeps the protections of
ADR-004: loopback only, exact \texttt{Host}, allowed \texttt{Origin} and a bearer
token on every request. VST3 plugins cannot create tracks; REAPER scripts can.

\textbf{Decision:}
\begin{itemize}
  \item \textbf{Discovery file.} When the engine starts it writes
        \texttt{\%LOCALAPPDATA\%\textbackslash Cover\_Maker\textbackslash run\textbackslash engine.json} with
        \texttt{schema}, \texttt{host} (always 127.0.0.1), \texttt{port},
        \texttt{pid}, \texttt{version} and a \emph{client token}. The file is
        written to a temporary name and renamed, so readers never see half a
        file, and deleted when the engine stops (only if it still names this
        engine's PID). A development build writes the same file under
        \texttt{Cover\_Maker.dev}; clients look there after the installed
        app's folder.
  \item \textbf{Separate client token.} The client token is 32 random bytes
        generated by the engine, distinct from the shell's session token, which
        still never touches the disk. The API accepts both; all other checks are
        unchanged, so a web page still cannot reach the engine (wrong
        \texttt{Host} or \texttt{Origin}) even if it guessed the port.
  \item \textbf{User-only folder.} The \texttt{run} folder's inherited permissions
        are removed and only the current user's SID is granted access
        (\texttt{icacls /inheritance:r}), so other accounts, including services
        running as other users, cannot read the token. The trust boundary is the
        Windows user account: any program that already runs as the user could
        also read the shell's memory.
  \item \textbf{Clients validate.} A client treats the file as stale unless
        \texttt{GET /health} with the token succeeds, and it sends the token in a
        header, never on a command line (the REAPER script passes it to
        \texttt{curl.exe} through a header file in the user-only folder, deleted
        after each call).
  \item \textbf{Opt-out.} Setting \texttt{COVER\_MAKER\_DAW\_ACCESS=0} disables
        the client token and the file.
  \item \textbf{REAPER} uses a Lua ReaScript: the selected item's file (or a
        typed file or link) goes to \texttt{POST /jobs}; the script polls the job
        with \texttt{reaper.defer} and inserts one track per stem below the
        source track, placed where the source item's start would be, in the
        order the stems were asked for.
  \item \textbf{VST3 plugin} (\texttt{daw/vst3}) is a client of the same API,
        written in Rust with NIH-plug and an egui editor: it separates a file
        or link, previews each stem through the track and lets the user drag
        the files into tracks (OLE drag and drop of a shell data object,
        \texttt{CF\_HDROP}), because the plugin standard cannot create tracks.
        It only connects to 127.0.0.1 and does not follow redirects.
  \item \textbf{Plugin licence.} NIH-plug is ISC, but its VST3 export uses the
        \texttt{vst3-sys} bindings, which are GPLv3. The product owner chose
        (2026-11-05) to ship the plugin as a separate GPL-3.0-or-later program
        and keep the rest of \CM{} under MIT. The alternative, a hand-written
        MIT wrapper over the \texttt{vst3} crate generated from the VST 3.8 SDK
        (MIT since October 2025), was rejected for its cost and unsafe code;
        it remains possible if NIH-plug moves to MIT bindings.
\end{itemize}

\textbf{Alternatives:} a fixed port (collides and is guessable); a named pipe
with an ACL (would need a second transport in the engine and every client); the
Windows Credential Manager (not reachable from Lua); passing the session token
to clients (one leak would expose the shell).

\textbf{Consequences:} third-party code running as the user can drive the
engine, which is acceptable because it could already run the engine itself;
a stale file after a crash is harmless because the port is closed or answers
401 to the old token; the protocol is plain HTTP, so other DAWs can be added
without engine changes.
